% !TEX root = ../../main.tex
%
% Exact Bernoulli witnesses and the open Virasoro residue comparison.

\providecommand{\Vir}{\mathrm{Vir}}
\providecommand{\ClaimStatusDefinitional}{}
\providecommand{\ClaimStatusProvedHere}{}
\providecommand{\ClaimStatusOpen}{}

\chapter{Bernoulli Witnesses and Virasoro Residue Comparison}
\label{ch:higher-kummer-arithmetic-duality}

\section*{Two arithmetic objects}

An odd prime \(p\) is Kummer-irregular when
\(p\mid\operatorname{num}(B_{2m})\) for an even index
\(2\le2m\le p-3\).  This classical condition applies directly to
the Bernoulli coefficients of the \(\mathrm{SL}_2\) Verlinde
polynomials.  A comparison with Virasoro higher residues begins after
the ordered-residue package
\(\mathsf H_{\mathrm{res}}(\Vir_c;X)\) constructs the scalar
coordinates and the arithmetic package fixes their primitive integer
normalization.

The Bernoulli window through \(B_{22}\) contains the set
\begin{equation}
 \mathrm{IrrKum}_{\le22}
 :=\{691,3617,43867,283,617,131,593\}.
 \label{eq:higher-irrkum-tower}
\end{equation}

\section{Exact Bernoulli witnesses}

\begin{theorem}[Kummer-irregular primes in the Bernoulli window]
\ClaimStatusProvedHere
\label{thm:higher-kummer-z-g-presence}
\emph{Type signature:}
\textup{(Open/modular quadrants, \(\mathrm{SL}_2\) Verlinde-polynomial
presentation, level \(5\),
\(\mathsf H_{Z_g}^{\mathrm{Bern}}\)).}
Here \(\mathsf H_{Z_g}^{\mathrm{Bern}}\) is the leading-coefficient
formula of Theorem~\ref{thm:z-g-leading-coefficient-bernoulli}
together with the standard Bernoulli normalization.
For every prime \(p\) in~\eqref{eq:higher-irrkum-tower}, the first
displayed Bernoulli witness and its Verlinde genus are:
\begin{center}
\begin{tabular}{rccc}
\hline
\(p\)&\(2m\)&\(g(p)=m+1\)&\(\lvert\operatorname{num}(B_{2m})\rvert\)\\
\hline
\(691\)&\(12\)&\(7\)&\(691\)\\
\(3617\)&\(16\)&\(9\)&\(3617\)\\
\(43867\)&\(18\)&\(10\)&\(43867\)\\
\(283\)&\(20\)&\(11\)&\(174611=283\cdot617\)\\
\(617\)&\(20\)&\(11\)&\(174611=283\cdot617\)\\
\(131\)&\(22\)&\(12\)&\(854513=11\cdot131\cdot593\)\\
\(593\)&\(22\)&\(12\)&\(854513=11\cdot131\cdot593\)\\
\hline
\end{tabular}
\end{center}
Each \(p\) divides the numerator of the leading coefficient of
\(Z_{g(p)}(k)\), written as a polynomial in \(n=k+2\).
\end{theorem}

\begin{proof}
The even Bernoulli numerators through \(B_{22}\) are
\[
\begin{array}{c|ccccccccccc}
2m&2&4&6&8&10&12&14&16&18&20&22\\ \hline
\lvert\operatorname{num}(B_{2m})\rvert
&1&1&1&1&5&691&7&3617&43867&174611&854513.
\end{array}
\]
The displayed factorizations give the seven primes and their first
indices in this window.  Every witness satisfies \(2m\le p-3\), so
Kummer's criterion applies.

By Theorem~\ref{thm:z-g-leading-coefficient-bernoulli}, the relevant
Verlinde coefficient equals
\[
 (-1)^g\frac{2^{g-1}B_{2g-2}}{(2g-2)!}.
\]
For the seven rows, \(p>2g-2\); hence the rational multiplier is a
\(p\)-adic unit.  The Bernoulli numerator factor therefore survives
in the reduced coefficient.
\end{proof}

\begin{remark}[The factor (11) and the Kummer window]
\label{rem:regular-factor-11}
The factor \(11\mid\operatorname{num}(B_{22})\) lies beyond the
Kummer index window for \(p=11\), since the defining range ends at
\(p-3=8\).  This illustrates the role of the index bound in the
cyclotomic criterion.
\end{remark}

\section{The Virasoro coefficient census}

\begin{definition}[Arithmetic residue package]
\ClaimStatusDefinitional
The package \(\mathsf H_{\mathrm{Kum}}^{\le13}\) consists of:
\begin{enumerate}[(i)]
\item \(\mathsf H_{\mathrm{res}}^{\le13}(\Vir_c;X)\) and its scalar
coordinates;
\item a proved rational-function presentation in every arity;
\item a canonical primitive integer normalization of each numerator;
\item exact factorization certificates from two independent integer
oracles;
\item compatibility with changes of residue representative and gauge.
\end{enumerate}
For a completed package, let \(C_r\) be the coefficient set of the
primitive numerator polynomial and let
\(\mathcal P_{\le13}:=\bigcup_{4\le r\le13}\operatorname{Supp}(C_r)\).
\end{definition}

\begin{problem}[Compute the Virasoro prime census through arity \(13\)]
\ClaimStatusOpen
\label{thm:higher-kummer-s-r-absence-through-r-13}
\emph{Type signature:}
\textup{(Open quadrant, ordered residue-bar and primitive numerator
presentations, level \(2\),
\(\mathsf H_{\mathrm{res}}^{\le13}
+\mathsf H_{\mathrm{Kum}}^{\le13}\)).}

Construct \(\mathsf H_{\mathrm{Kum}}^{\le13}\), compute the sets
\(C_r\), and decide the finite separation condition
\begin{equation}
 \mathcal P_{\le13}\cap\mathrm{IrrKum}_{\le22}=\varnothing.
 \label{eq:higher-kummer-s-r-absence}
 \label{eq:char-s-r-disjoint}
\end{equation}
The computation includes chain-level residue representatives and
factorization certificates for every primitive coefficient.
\end{problem}

\begin{remark}[Characteristic-prime certificate]
\label{rem:characteristic-s-r-primes}
The characteristic-prime set is an output of the arithmetic residue
package.  Its certificate records each coefficient, its primitive
normalization, and its exact factorization.
\end{remark}

\begin{remark}[First overlap boundary]
\label{rem:2111-sharpness}
After the coefficient census is constructed, the first overlap with
the full set of Kummer-irregular primes defines a further arithmetic
invariant of the residue tower.  This boundary is computed from the
certified sets \(C_r\).
\end{remark}

\section{Finite Bernoulli--Virasoro comparison}

\begin{problem}[Determine the finite genus--residue separation]
\ClaimStatusOpen
\label{thm:higher-kummer-refined-duality}
\emph{Type signature:}
\textup{(Open/modular quadrants, primitive residue numerator and
Bernoulli presentations, levels \(2\leftrightarrow5\),
\(\mathsf H_{\mathrm{Kum}}^{\le13}
+\mathsf H_{Z_g}^{\mathrm{Bern}}\)).}
For \(p\in\mathrm{IrrKum}_{\le22}\), define
\[
 r_0(p):=\inf\{r\ge4:p\in\operatorname{Supp}(C_r)\},
 \qquad
 g(p):=\min\{g\ge2:p\mid\operatorname{num}(B_{2g-2})\}.
\]
Use the convention \(\inf\varnothing=\infty\).
Theorem~\ref{thm:higher-kummer-z-g-presence} gives
\[
 (g(691),g(3617),g(43867),g(283),g(617),g(131),g(593))
 =(7,9,10,11,11,12,12).
\]
Complete the Virasoro census and determine whether it yields
\begin{equation}
 g(p)\le12<14\le r_0(p)
 \qquad(p\in\mathrm{IrrKum}_{\le22}).
 \label{eq:boundary-inequality}
\end{equation}
\end{problem}

\begin{remark}[Weight and arity boundary]
\label{rem:higher-kummer-weight-boundary}
The Bernoulli index \(2g-2\) and the residue arity~\(r\) are separate
gradings.  A comparison theorem records the map between them together
with its normalization; the finite inequality above is the first
test case.
\end{remark}
